% !TEX root = IE3_expI_report.tex
\documentclass[lualatex,a4paper,10pt]{jlreq}
% Shared LuaLaTeX preamble.
\usepackage{luatexja-fontspec}
\setmainfont{TeX Gyre Termes}
\setsansfont{TeX Gyre Heros}
\setmonofont{Latin Modern Mono}
\setmainjfont[BoldFont=HaranoAjiGothic-Medium]{HaranoAjiMincho-Regular}
\setsansjfont[BoldFont=HaranoAjiGothic-Bold]{HaranoAjiGothic-Medium}
\usepackage[top=22mm,bottom=22mm,left=23mm,right=23mm]{geometry}
\usepackage{amsmath,amssymb,booktabs,array,longtable,tabularx}
\usepackage{graphicx,xcolor,tikz,circuitikz}
\usetikzlibrary{arrows.meta,positioning,calc}
\usepackage{enumitem,fancyhdr,listings,needspace}
\usepackage[unicode,colorlinks=true,linkcolor=labblue,urlcolor=labteal]{hyperref}
\usepackage{bookmark}
\definecolor{labblue}{RGB}{30,70,112}
\definecolor{labteal}{RGB}{0,100,105}
\definecolor{labgray}{gray}{0.95}
\ModifyHeading{section}{font={\large\sffamily\bfseries\color{labblue}}}
\ModifyHeading{subsection}{font={\normalsize\sffamily\bfseries\color{labteal}}}
\setlength{\parindent}{1\zw}
\setlength{\parskip}{3pt}
\setlength{\emergencystretch}{2em}
\renewcommand{\arraystretch}{1.35}
\setlist{itemsep=3pt,topsep=4pt,leftmargin=2\zw}
\newcolumntype{Y}{>{\raggedright\arraybackslash}X}
\newcolumntype{C}[1]{>{\centering\arraybackslash}p{#1}}
\pagestyle{fancy}
\fancyhf{}
\fancyhead[L]{\small\color{labblue} IE3 後期・電子工学実験}
\fancyhead[R]{\small テーマ\LabCode}
\fancyfoot[C]{\thepage}
\setlength{\headheight}{16pt}
\DeclareOldFontCommand{\rm}{\normalfont\rmfamily}{\mathrm}
\lstset{basicstyle=\small\ttfamily,columns=fullflexible,keepspaces=true,
 breaklines=true,frame=single,backgroundcolor=\color{labgray},
 numbers=left,numberstyle=\tiny,showstringspaces=false,aboveskip=8pt,belowskip=8pt}
\newcommand{\blank}{\rule{22mm}{0.3pt}}
\newcommand{\answerline}{\par\noindent\rule{\linewidth}{0.25pt}\par}
\newcommand{\writing}[1]{\par\Needspace{\dimexpr#1+5mm\relax}\noindent%
 \fcolorbox{labblue!35}{white}{\parbox[c][#1][t]{\dimexpr\linewidth-2\fboxsep-2\fboxrule\relax}{\vspace{2mm}\noindent\strut\hfill\par}}\par}
\newcommand{\hint}[1]{\par\smallskip\noindent{\sffamily\bfseries\color{labteal}記述の視点：}#1\par}
\newcommand{\note}[1]{\par\smallskip\noindent\fcolorbox{labblue!45}{labblue!5}{%
 \parbox{\dimexpr\linewidth-2\fboxsep-2\fboxrule\relax}{#1}}\par\smallskip}
\newcommand{\eqerror}{\varepsilon=\frac{x_{\mathrm{meas}}-x_{\mathrm{th}}}{x_{\mathrm{th}}}\times100\ \%}
\newcommand{\LabSource}{徳山工業高等専門学校 情報電子工学科，
 『2026年度 電子工学実験（3年後期）テキスト』，テーマ\LabCode，
 \expandafter\path\expandafter{\SourcePdf}。}
\newcommand{\reporttitle}{
 \noindent{\small\sffamily\color{labteal}実験報告書・記入ガイド}\par
 \noindent{\LARGE\sffamily\bfseries \LabCode\quad\LabTitle}\par\smallskip
 \noindent 氏名：\blank\quad 班：\blank\quad 実験日：\blank\par
 \note{目的・原理・手順を確認し、実測値と自分の考察を記入する。
 考察のガイドは、検討する事象・使う測定結果・記述の方針を示す。}
}
\newcommand{\prelabtitle}{
 \noindent{\small\sffamily\color{labteal}事前学習・前提知識プリント}\par
 \noindent{\LARGE\sffamily\bfseries \LabCode\quad\LabTitle}\par\smallskip
 \noindent 氏名：\blank\quad 班：\blank\quad 予習日：\blank\par
 \note{用語、回路の動き、計算、測定表への記入の順に読む。
 例題の値は説明用であり、実験結果ではない。}
}
\newcommand{\tocrow}[3]{\hyperref[#1]{#2}& #3 &\pageref*{#1}\\[2mm]}
\newcommand{\documentmap}[1]{
 \section*{目次と概要}
 \noindent 青色の項目をクリックすると該当箇所へ移動する。\par
 {\small\begin{longtable}{@{}p{0.29\linewidth}p{0.61\linewidth}r@{}}
 \toprule {\color{labblue}\bfseries 項目} & {\color{labblue}\bfseries 読むと分かること} & 頁\\\midrule
 \endhead #1\bottomrule
 \end{longtable}}
 \clearpage
}
\newenvironment{terms}{\par\smallskip\begin{longtable}{@{}p{0.23\linewidth}p{0.73\linewidth}@{}}
 \toprule {\color{labteal}\bfseries 用語}&{\color{labteal}\bfseries 意味・回路での役割}\\\midrule\endhead}
 {\bottomrule\end{longtable}}
\newcommand{\term}[2]{\textbf{#1}&#2\\[2mm]}
\newcommand{\guidepoint}[4]{
 \Needspace{32mm}\subsection{#1}
 \noindent{\bfseries\color{labteal}考える事象：}#2\par
 \noindent{\bfseries\color{labteal}参照する結果：}#3\par
 \noindent{\bfseries\color{labteal}記述の方針：}#4\par
}
\newcommand{\reportclosing}[1]{
 \clearpage\section{結論のまとめ方}\label{sec:closing}
 \hint{#1}
 \ConclusionGuide
 \begin{enumerate}
 \item 目的に対応する最も重要な実測結果を、測定条件と数値を添えて述べる。
 \item 原理と一致した点・一致しなかった点を示す。原因は確認済みの事実と推測を分ける。
 \item 実施範囲で言えることと、未確認の条件を一文ずつまとめる。
 \end{enumerate}
 \note{書き出しの型：「○○の条件で△△を測定したところ、□□となった。
 これは◇◇と整合する。ただし××は確認しておらず、一般化には追加測定が必要である。」
 空欄は自分の実測値と根拠で埋める。}
 \noindent 結論（実験後に記入）：\writing{30mm}
 \noindent 改善案・次に確認したい条件：\writing{16mm}
 \section*{参考文献}\label{sec:references}
 \begin{enumerate}
 \item \LabSource
 \item 使用したデータシート・書籍・ウェブ資料（著者、題名、該当箇所、URL、参照日）を追加する。
 \end{enumerate}
}
\newcommand{\prelabclosing}{
 \section*{準備・出典}\label{sec:prelabsource}
 \begin{itemize}[nosep]
 \item 資料、筆記具、記録用紙、電卓と、実験条件・機器設定の記録欄。
 \item 確認問題の解答と、分からない用語・回路点への具体的な質問。
 \end{itemize}
 {\small\noindent\LabSource\par}
}
\newcommand{\simpleflow}[1]{
 \begin{center}
 \begin{tikzpicture}[node distance=5mm and 5mm,
 every node/.style={font=\small},box/.style={draw=labblue,fill=labblue!4,rounded corners=1mm,
 align=center,minimum height=10mm,text width=30mm}]
 #1
 \end{tikzpicture}
 \end{center}
}

\newcommand{\LabCode}{I}
\newcommand{\LabTitle}{A/D変換とD/A変換}
\newcommand{\SourcePdf}{IE3_expI_A4.pdf}
\newcommand{\ConclusionGuide}{分解能は符号と電圧の表、標本化は周波数を変えた波形、保持は制御パルスとの対応を根拠にまとめる。
理想スケールと実測の差、再構成の条件、未確認の誤差要因をそれぞれ区別して述べる。\par}
\begin{document}
\reporttitle
\documentmap{
\tocrow{sec:auto1}{1 目的}{量子化・標本化・重み付けによる変換を実験で確かめる。}
\tocrow{sec:auto2}{2 原理}{量子化・標本化・重み付けによる変換の仕組みと成立条件を整理する。}
\tocrow{sec:auto3}{3 装置・条件}{機器・定数・測定点を確認し、資料の順序で操作する。}
\tocrow{sec:auto4}{4 方法}{機器・定数・測定点を確認し、資料の順序で操作する。}
\tocrow{sec:auto5}{5 結果}{符号と電圧、折返し、保持波形を表や図に記録する。}
\tocrow{sec:auto6}{6 考察：結果を根拠に説明する}{符号と電圧、折返し、保持波形を根拠に、比較と原因の検討を進める。}
\tocrow{sec:closing}{7 結論のまとめ方}{主要な結果、成立範囲、未確認の条件を短くまとめる。}
\tocrow{sec:references}{参考文献}{使用した資料と外部の根拠を示す。}
}

\section{目的}\label{sec:auto1}
3ビットと8ビットのA/D・D/A変換を観測し、量子化、重み付け、サンプリング、サンプルホールドの動作を理解する。
\section{原理}\label{sec:auto2}
nビットの理想D/A変換は、フルスケール$V_{\rm FS}$と符号Dに対し
\[
 V_o=V_{\rm FS}\frac D{2^n},\quad
 \mathrm{LSB}=\frac{V_{\rm FS}}{2^n},\quad 0\le D\le2^n-1
\]
となる。R--2Rラダーでは各ビットの寄与が1/2、1/4、…と重み付けされる。
3ビットの資料回路は比較器とエンコーダによる並列比較形である。
\note{資料中のADC0804の方式説明は補足が必要である。メーカー資料では8ビット逐次比較形ADCとされている。3ビットの並列比較回路と同一方式ではない。}
帯域制限された信号の理想的な再構成には$f_s>2f_{\max}$が必要であり、これ以下では折返しが起こり得る。
\section{装置・条件}\label{sec:auto3}
3ビット回路：LM339、74LS148、R--2R（10 k$\Omega$/20 k$\Omega$）、LM358。
8ビット回路：資料指定のADC/DACボード、9 V供給と内部5 V、基準2.5 V。DAC出力増幅を含む理想スケールは5 V。
信号発振器、オシロスコープ、電圧計、実際の基準・電源値を記録する。
\writing{16mm}
\section{方法}\label{sec:auto4}
\begin{enumerate}
\item 3ビットADCの入力電圧と符号を調べ、組み立てたR--2R DACの各符号の出力を測る。入力100 Hz、5 Vpp、オフセット2.5 Vの正弦波でも観測する。
\item 8ビット回路で50 Hz、5 Vpp、オフセット2.5 Vのクロックを与え、指定符号とDAC出力を確認する。
\item 入力50 Hzの正弦波、サンプリング500〜2500 Hzで変換波形を比較する。出力を基準としたトリガやSTOPを使い観測条件をそろえる。
\item 110 Hz入力／80 Hzサンプリングを5周期程度記録する。50 Hz入力／60 Hzサンプリングも観測する。
\item サンプルホールドで100 Hz正弦波と、50 Hz・デューティ20\%のパルスを用い、入力、パルス、保持出力を同じ時間軸で比較する。
\end{enumerate}
\clearpage
\section{結果}\label{sec:auto5}
表1：3ビット変換。ADCの入力範囲とDAC出力を区別する。
\par\noindent\begin{tabularx}{\linewidth}{YYYY}
\toprule 符号D & ADC入力範囲/V & DAC理論/V & DAC実測/V\\\midrule
0&&0&\\1&&0.625&\\2&&1.250&\\3&&1.875&\\
4&&2.500&\\5&&3.125&\\6&&3.750&\\7&&4.375&\\\bottomrule
\end{tabularx}
表2：8ビット、5 Vスケールの理論と測定。
\par\noindent\begin{tabularx}{\linewidth}{YYYY}
\toprule 符号/16進 & 理論出力/V & 実測入力/V & 実測出力/V\\\midrule
00&0&&\\01&0.01953&&\\02&0.03906&&\\04&0.07813&&\\
08&0.15625&&\\10&0.31250&&\\20&0.62500&&\\
40&1.25000&&\\80&2.50000&&\\FF&4.98047&&\\\bottomrule
\end{tabularx}
図1：3ビットおよび8ビットの入出力波形。
\writing{38mm}
\clearpage
表3：入力50 Hzに対するサンプリング比較（観測した条件を追記する）。
\par\noindent\begin{tabularx}{\linewidth}{YYYY}
\toprule $f_s$/Hz & 1周期の標本数 & 出力波形の特徴 & 記録図\\\midrule
500&&&\\1000&&&\\1500&&&\\2000&&&\\2500&&&\\60&&&\\\bottomrule
\end{tabularx}
図2：110 Hz／80 Hz、および50 Hz／60 Hzの折返し波形。
\writing{38mm}
図3：サンプルホールドの入力、パルス、出力（同一軸）。
\writing{35mm}
\clearpage\section{考察：結果を根拠に説明する}\label{sec:auto6}
\guidepoint{符号と出力電圧・LSB}
{3ビットと8ビットの刻み、最大値、理論との差はどうなるか。}
{表1・2の各符号とDAC実測、零符号、基準電圧、出力増幅率。}
{符号を十進へ直して理論を計算する。隣接点や1ビット符号から刻みを調べ、オフセットとスケールの違いを候補別に検討する。}
\guidepoint{ADCの段階とR--2Rの重み}
{入力の境界と各ビットの寄与が説明できるか。}
{ADC入力範囲、3ビット出力、R/2Rの実値、比較器・エンコーダの接続。}
{各ビットだけを有効にした寄与を、重ねの理と端子等価から導く。生の負論理端子と表示符号を区別し、最大電圧だけで全段の精度を評価しない。}
\guidepoint{サンプリング周波数と折返し}
{段階波形の見やすさと、入力とは別の周波数に見える現象がどう違うか。}
{50 Hz入力の500〜2500 Hz系列、110/80 Hzと50/60 Hzの波形、実際の周期。}
{一周期の標本数と折返し周波数を計算して照合する。標本化定理の帯域条件と、実回路の変換時間・フィルタの制約を分ける。}
\guidepoint{サンプルホールドの取得と保持}
{入力へ追従する区間と保持する区間は制御パルスに対応するか。}
{同一時間軸の三波形、取得極性、周期・デューティ、保持中の電圧変化。}
{保持開始時の値を入力に対応させる。同位相取得やトリガ条件を確認し、垂れ下がり等を評価した場合は読取り幅と保持時間も示す。}
\subsection{資料の課題・追加検討}
\begin{enumerate}
\item ADCの階段状特性と、3ビット・8ビットのLSBを説明する。最大符号が5 Vちょうどに達しない理由も示す。
\item 重ねの理とテブナン等価を用いてR--2Rラダー各ビットの寄与を導く。
\item 110 Hz／80 Hz、50 Hz／60 Hzの折返しを計算し、観測と比較する。
\item 1周期10標本等の見やすさと、標本化定理による再構成の条件を区別する。
\item サンプル区間と保持区間を示し、保持中の変化や誤差を検討する。
\end{enumerate}
\reportclosing{分解能、サンプリング周波数、保持動作による波形の変化をまとめる。}
\end{document}
